
\فصل{جدایی‌پذیری مخلوط‌های غیرخطی}
\label{seprability}
در این بخش به بررسی قضایای مربوط به جدایی پذیری مخلوط‌های غیرخطی می‌پردازیم.

\section{
	روش تمایز بین زمان‌ها
\lr{(TCL)}
}
روش تمایز بین زمان‌ها یا 
\lr{Time-Contrastive Learning (TCL)} 
روشی است برمبنای آموزش یک استخراج‌گر ویژگی\LTRfootnote{Feature Extractor} با استفاده از یک طبقه‌بند\LTRfootnote{Classifier}
\lr{Multinomial Logistic Regression}
که می‌خواهد بین قطعه‌های زمانی سیگنال زمانی، تفکیک قائل‌شود. این روش در 
\cite{TCL}
معرفی شده‌است. در این بخش، در ابتدا به معرفی 
\lr{TCL}
پرداخته و سپس نشان می‌دهیم این روش می‌تواند دسته‌ای از مسائل جداسازی کور مخلوط‌های غیرخطی را حل نماید.

\subsection{الگوریتم 
TCL
}

الگوریتم زیر را در نظر بگیرید:

\begin{enumerate}%\def\labelenumi{A\theenumi)}
	\setlength{\itemsep}{-0.5mm}%\setlength{\paraindent}{0mm}
	\item 
	سیگنال‌های زمانی 
	$\x(t)$
	را به تعدادی سگمنت‌ زمانی تقسیم کنید. (پنجره‌های زمانی که اندیس آن‌ها $\tau=1,\ldots,T$) است. طول این پنجره‌ها را برابر در نظر بگیرید. فرض کنید در هر پنجره‌ی زمانی، $L$ استپ زمانی داریم.
	\item 
مقادیری از سیگنال 
$\x(t)$
که در بازه‌ی زمانی 
$\tau$
قرار دارد، برچسب 
$\tau$
بدهید. اگر در هر سگمنت $L$ استپ زمانی داشته باشیم، برای هر برچسب در دیتاست تولید شده، $L$ نمونه داریم. هر یک از این نمونه‌ها یک بردار $n$ بعدی است ($\x(i)$).

\item
یک استخراج‌گر ویژگی 
$\mathbf{h}(\x_t;\thetab)$ 
و یک 
\lr{MLR}
را برای طبقه‌بندی دیتاستی که در مرحله‌ی تولید کرده‌اید آموزش دهید. برای 
\lr{MLR}
از تابع رگرسیون 
$ w_\tau^T \h(\x_t; \thetab)+b_\tau$
استفاده نمایید.

\end{enumerate}

در این جا هدف استخراج‌گر ویژگی آن است که ویژگی‌هایی از داده بیابید که به 
\lr{MLR}
کمک کند طبقه‌بندی بهتری انجام دهد.  شهوداً می توان انتظار داشت که این استخراج‌گر باید نمایش مناسبی از سیگنال و رابط زمانی موجود در آن ر ا یاد بگیرد. در عمل، از آن‌جا که انتظار داریم تابع 
$\h$
توانایی تقریب هر تابع غیرخطی را داشته باشد، از یک 
\lr{MLP}
با حذف لایه‌ی آخر آن، به عنوان $\h$ استفاده می‌کنیم. در بخش بعد اثبات می‌‌کنیم که این استخراج‌گر ویژگی، در عمل، جداسازی کور منابع را انجام داده است. شکل 
\eqref{TCL}،
روند کلی این روش را نمایش می‌دهد.
\begin{figure}
\centering
\includegraphics[scale=1.1]{TCL.png}
\caption{شمای کلی روش تمایز بین زمان‌ها
(TCL)
}
\label{TCL}
\end{figure}

\subsection{
\lr{TCL}
به عنوان تقریبی از نسبت
\lr{log-pdf}
ها
}

در یک 
\lr{MLR}،
احتمال
\lr{posterior}
برای اینکه دیتای 
$\x_t$
عضو دسته‌ی 
$\tau$
باشد، برابر است با:
\begin{align}
%    \ell(\thetab, \m W, \ve b) = 
p(C_t=\tau | \x_t; \thetab,  \mathbf {W},  \mathbf {b})= 
%\nonumber \\
\frac{\exp( w_\tau^T \h(\x_t; \thetab)+b_\tau)}
{1+\sum_{j=2}^{T} \exp( w_j^T \h(\x_t; \thetab)+b_j)}  \label{eq:ell2}
\end{align}
که در آن 
$C_t$
برچسب دیتا در زمان 
$t$،
$\x_t \in \R^n$
بردار دیتا در زمان $t$،
$\theta$
پارامتر‌های استراج‌گر ویژگی،
$\mathbf{W} = [\mathbf{w}_1, \mathbf{w}_2, \dots, \mathbf{w}_T] \in 
\R^{m \times T}$
وزن‌ها 
و 
$\mathbf{b} = [b_1, \dots b_T]^\top$
بایاس‌های 
\lr{MLR}
هستند. از طرف دیگر، با قانون بیز، می‌توان نوشت:
\begin{equation}
p(C_t=\tau | \x_t) = \frac{p_\tau(\x_t)p(C_t = \tau)}
{\sum_{j=1}^{T} p_j(\x_t)p(C_t = j)},
\label{eq:pc_i}
\end{equation}
که در آن 
$p(C_t = \tau)$
توزیع پیشین برچسب یک سگمنت است و 
$p_\tau(\x_t) = p(\x_t| C_t = \tau)$
است. اگر استخراج‌گر ویژگی، ویژگی
\lr{Universal Approximation Capacity}
را داشته و بیشمار داده در اختیار باشد، 
\eqref{eq:ell2}
به 
\eqref{eq:pc_i}
همگرا می‌شود. با مساوی قرار دادن این دو عبارت، داریم:
\begin{equation} 
w_\tau^T \h(\x_t; \thetab) + b_\tau 
= \log p_\tau(\x_t) - \log p_1(\x_t) + \log \frac{p(C_t = \tau)}{p(C_t = 1)} ,
\label{eq:relation}
\end{equation}

اگر طول سگمنت‌ها برابر باشد، عبارت سومین جمله‌ی سمت راست تساوی نیز برابر با صفر می‌شود، در نتیجه:
\begin{equation} 
w_\tau^T \h(\x_t; \thetab) + b_\tau 
= \log p_\tau(\x_t) - \log p_1(\x_t)
\label{eq:relation2}
\end{equation}

\subsection{	استفاده از
\lr{TCL}
	 در جداسازی کور منابع}
مجدداً مدل مطرح شده در فصل
\eqref{probdef}
را در نظر بگیرید:
\begin{equation}
\s_t = \big(s_1(t), s_2(t), \dots, s_n(t)\big);\;\;\;\;
\x_t = \mathbf{f}(\s_t)
\end{equation}
 فرض‌کنید سیگنال‌های 
$s_i(t)$
غیر ایستان باشند.  همچنین 
$s_i(t)$
ها را فرآیندهایی مشترکاً مستقل در نظر بگیرید. فرض کنید، غیر ایستان بودن فرآیند‌ها آنقدر کند است که می‌توانیم توزیع را در هر سگمنت از دیتا ثابت بگیریم. 
\lr{log-pdf}
منبع با اندیس 
$i$
در سگمنت
$\tau$
را به صورت زیر تعریف می کنیم:
\begin{equation} 
\log p_\tau(s_i)=  q_{i,0}(s_i) + \sum_{v=1}^V \lambda_{i,v}(\tau) q_{i,v}(s_i) - \log Z( \lambda_{i,1}(\tau),\ldots,\lambda_{i,v}(\tau))
\label{eq:p_tau}
\end{equation}
اندیس $t$ برای سادگی حذف شده‌است.
در این عبارت،
$q_{i, 0}$
یک
\lr{stationary baseline}
برای توزیع منابع است. توابع 
$q_{i, v}$
به ازای 
$ v\geq 1$
توابع غیرخطی اسکالری هستند که تابع 
$\tau$
بوده و مدل‌کننده‌ی غیر ایستان بودن منبع هستند.
$Z$
نقش نرمالایز کننده‌ی توزیع احتمال را بر عهده دارد که در تمامی اثبات‌ها حذف می‌شود و بی‌اثر است.

\begin{example}
برای مثال، فرض کنید 
$q_{i, 0} = 0$
است و $V$ را برابر یک در نظر بگیرید. تابع
$q_{i, 1} (s_i) = -\frac{s_i^2}{2}$
مدل کننده‌ی یک منبع گوسی است که واریانس آن در حال تغییر است. اگر 
$q_{i, 1} (s_i) = -|s_i|$
در نظر می‌گرفتیم، منبع یک منبع لاپلاسی می‌شد.
\end{example}

حال می‌خواهیم به بررسی روش استفاده از 
\lr{TCL}
برای حل جداسازی کور بپردازیم.

\begin{theorem}
در مدل 
\eqref{eq:p_tau}
فرض کنید 
$V = 1$
است،
$q_{i, 0} = 0$
و
$q \triangleq q_{i, 1}$،
یعنی منابع به صورت زیر تولید شده‌اند:
\begin{equation}
\log p_\tau(s_i)=    \lambda_{i,1}(\tau) q(s_i) - \log Z( \lambda_{i,1}(\tau))
\label{eq:p_tau_simple}
\end{equation}
\begin{itemize}
\item
فرض کنید 
\lr{TCL}
را بر  روی دیتای تولیدشده از مدل فوق اعمال کنیم و در آن، بُعد استخراج‌گر ویژگی، 
$\h$
را برابر تعداد منابع (که برای تعداد سری‌های زمانی مشاهده‌شده است) در نظر بگیریم. 
\item
فرض کنید ماتریس 
$\mathbf{L}$
با المان‌های 
$$[\mathbf{L}]_{\tau, i} = 
\lambda_{i, 1}(\tau) - \lambda_{i, 1}(1) \
\;\;\;\;
\tau = 1, \dots, T, \;\;
i = 1, \dots, n
$$
داری رتبه‌ی ستونی کامل باشد\footnote{
	توجه کنید که از آن‌جایی که ماتریس 
$\mathbf{L}$
یک ماتریس بلند است و حتی با فرض مذکور، سگمنت‌های زمانی زیادی می‌توانند توزیع یکسانی داشته باشند.}.
\end{itemize}
در این حالت، در حدی که بیشمار نمونه‌ی زمانی از فرآیند های تصادفی در اختیار داشته داریم، ماتریس 
$\mathbf{A} \in \R^{n\times n}$
و
$\mathbf{d} \in \R^n$
وجود دارند که:
\begin{equation}
q(\s_t) = \mathbf{Ah}(\x_t, \theta) + \mathbf{d}
\end{equation}
\end{theorem}

\begin{note}
شهود کامل بودن رتبه‌ی ماتریس 
$\mathbf{L}$
این است که واریانس منابع، مستقل از یکدیگر در حال تغییر هستند.
\end{note}
\begin{proof}
در ابتدا، توزیع احتمال 
$\x_t$
را بر حسب توزیع احتمال
$\s_t$
به دست می‌آوریم:
\begin{equation}
\log p_\tau(\x) = \sum_{i=1}^n \lambda_{\tau,i}q(g_i(\x)) 
+ \log | \det  \mathbf{J} \mathbf{g} (\x) | - \log Z(\lambda_\tau),
\label{eq:p_x}
\end{equation}
در این عبارت، برای سادگی از اندیس 
$t$ 
برای
$\x$
چشم پوشی کرده‌ایم. همچنین
\begin{equation}
\mathbf{g}(\x) = \big(g_1(\x), g_2(\x), \dots, g_n(\x)\big)^\top
\end{equation}
تابع وارون مخلوط‌کننده‌ی
$\mathbf{f}$
است (یعنی بنا بر تعریف
$s_i = g_i(\x)$).
ماتریس 
$\mathbf{Jg}$
نیز ژاکوبین تبدیل 
$\mathbf{g}$
است. با استفاده از 
\eqref{eq:relation}
می‌توان نوشت:
\begin{equation} 
\log p_\tau(\x)= \sum_{i=1}^n w_{\tau,i} h_i(\x) + b_\tau + \log p_1(\x) - c_\tau,
\label{eq:optlrpdfbias}
\end{equation}
در این معادله، 
$w_{\tau, i}$
و
$h_i{\x}$،
المان‌های 
$i$ام
بردار‌های 
$\mathbf{w}_\tau$
و
$\mathbf{h}(\x)$
هستند. در این عبارت، پارامتر‌های
$\theta$
در 
$\mathbf{h}$
را برای سادگی حذف کرده‌ایم. 
$c_\tau$
نیز عبارت آخر در سمت راست
\eqref{eq:relation}
است. با قرار دادن 
$\tau = 1$
در 
\eqref{eq:p_x}
داریم:
\begin{equation}
\log p_1(\x)= \sum_{i=1}^n \lambda_{1,i} q(g_i(\x)) + \log |\det \mathbf{ J g}(\x)| - \log Z(\lambda_1).
\label{eq:p1_x}
\end{equation} 
با جای‌گذاری این عبارت در 
\eqref{eq:optlrpdfbias}
داریم:
\begin{align} \label{eq:optlrpdf2}
\log p_\tau(\x) = \sum_{i=1}^n \left[w_{\tau,i} h_i(\x) + \lambda_{1,i} q(g_i(\x)) \right] \\
+ \log |\det\mathbf{ J g}(\x)| - \log Z(\lambda_1)  + b_\tau - c_\tau.
\end{align}

با مساوی قرار دادن 
\eqref{eq:optlrpdf2}
و
\eqref{eq:p_x}
به ازای هر 
$\tau$
داریم:
\begin{equation}
\sum_{i=1}^n \tilde{\lambda}_{\tau,i} q(g_i(\x)) 
= \sum_{i=1}^n w_{\tau,i} h_i(\x) + \beta_\tau,
\label{eq:maineq}
\end{equation}
در این معادله داریم
$\tilde{\lambda}_{\tau,i}= \lambda_{\tau,i}- \lambda_{1,i}$ 
و
 $\beta_\tau= \log Z(\lambda_\tau)-\log Z(\lambda_1)+b_\tau-c_\tau$.
با نوشتن این معادله‌ها برای 
$\tau$
های مختلف  به فرم ماتریس داریم:
\begin{equation}
\mathbf{ L} q(\s) %\begin{pmatrix} q(s_1)\\ q(s_2)\\ \vdots\\ q(s_n) \end{pmatrix} 
= \mathbf{W} \h(\x) + \mathbf{\beta}
\end{equation}
در این معادله،
$\mathbf{ L}$
همان ماتریس صورت قضیه است،
$\mathbf{ W}$
فرم ماتریسی
$w_{\tau, i}$
و
$\mathbf{\beta}$
بردار 
$\beta_\tau$
هاست.

با توجه به این‌که ماتریس
$\mathbf{ L}$
دارای رتبه‌ی ستونی کامل است، برای شبه-وارون آن داریم
$\mathbf{L}^+\mathbf{L} = \mathbf{I}$. 
با ضرب شبه-وارون در عبارت آخر به 
\begin{equation}
q(\s)
%\begin{pmatrix} q(s_1)\\ q(s_2)\\ \vdots\\ q(s_n) \end{pmatrix} 
= [\mathbf{L}^+  \mathbf{W}] \h(\x) + \mathbf{L}^+ \beta.
\end{equation}
می‌رسیم. که معادل گزاره‌ی صورت سوال است که در آن 
$\mathbf{A} = \mathbf{L}^+  \mathbf{W}$
یک ماتریس وارون‌پذیر است. اگر رتبه‌ی این ماتریس کامل نبود، توابع 
$q(s_i)$
وابسته‌ی خطی می‌بودند، که غیرممکن است زیرا تابع متغیر‌های 
$s_i$
هستند.
\end{proof}

\begin{corollary}
فرض‌ کنید
$\lambda_{i, 1}$
ها به صورت تصادفی تولید شده‌باشند. جدا کردن منابع 
$s_i$
می‌تواند با اعمال 
\lr{TCL}
و سپس اعمال یک الگوریتم خطی جداسازی کور بر روی خروجی‌های استخراج‌گر ویژگی
$\mathbf{h}$
ا نجام بگیرد، زیر $q$ یک کران بالا دارد (به خاطر انتگرال‌پذیر بودن) و در نتیجه 
$q(s_i)$
غیر گوسی است.
\label{corr}
\end{corollary}

\begin{note}
در حالت کلی، تابع $q$ وارون‌پذیر نیست و از نتیجه‌ی
\eqref{corr}
  نمی‌توان برای جداسازی 
$s_i$
ها استفاده کرد.
\end{note}
\begin{corollary}
فرض کنید 
$q(s)$
تابعی اکیداً یکنوا از 
$|s|$
باشد. در این صورت،  $s_i$ ها «قابل بازسازی» هستند.
\end{corollary}
\begin{proof}
برای انتگرال‌پذیر بودن
$p_\tau(s)$،
باید در حد
$|s| \to \infty$،
تابع
$q(s) \to -\infty$
برود.  در ضمن،
$q(s)$
باید یک ماکزیمم محدود داشته باشد. بدون کم شدن از کلیت مسئله، این ماکزیمم را برابر صفر در نظر بگیرید. برای یک $i$ مشخص، منیفلد
$q(g_i(\x)) = 0$
را در نظر بگیرید. این منیفلد، فضای 
$\x$
را به یک بخش با 
$\tilde{s_i} = q(s_i)$
و یک بخش با 
$\tilde{s_i} = -q(s_i)$
تقسیم می‌کند. با این 
$\tilde{s_i}$،
ما منابع را در حد یک تابع یکنوا بازسازی کرده‌ایم:
$\tilde{s_i} = c \;\mathrm{sign}(s_i)q(s_i)$
که در آن 
$c \in \{-1, +1\}$.
\end{proof}
\subsubsection{تعمیم برای منابع دیگر}
تعمیم این قضیه به منابعی که دارای یک 
\lr{stationary baseline}
ناصفر هستند، ساده است. تعمیم آن به حالت 
$V > 1$
ساده نیست و در پروژه‌ی کارشناسی ۲ مورد بررسی قرار خواهد گرفت.

\section{استفاده از متغیر‌های کمکی}

ایده‌ی اصلی این روش‌ها این است که اگر 
$n$
منبع وجود داشته باشند، $m$ متغیر 
$\u$
را به نحوی انتخاب می‌کنند که توابع 
$q_i$
وجود داشته باشند به نحوی که:
\begin{equation}
\log p(\s|\ub) = \sum_{i = 1}^{n} q_i(s_i, \ub)
\label{genprior}
\end{equation}
این ایده‌ می‌تواند تعمیم ایده‌ی 
\lr{(TCL)}
در نظر گرفته شود که در آن از شماره‌ی سگمنت به عنوان متعیر 
$\ub$
استفاده می‌شد. ایده‌ این است که مجدداً یک دیتاست بسازیم:
\begin{align}
\xaug=(\x,\uu)  \  \text{ vs. } \
\xaug^*=(\x,\uu^*) \label{xtildestar}
\end{align}
در این‌جا 
$\uu^*$
هم‌توزیع با 
$\uu$
است، اما از 
$\x$
مستقل است. در عمل 
$\uu^*$
را با جایگشت دادن نمونه‌های 
$\uu$
می‌سازیم. یک
\lr{Logistic Regression}
را با تابع رگرسیون
\begin{equation} \label{reggen}
r(\x,\uu)= \sum_{i=1}^n \psi_i(h_i(\x),\uu)
\end{equation}
یاد می‌گیریم که به ما احتمال این‌که یک بردار 
$\x$
و 
$\uu$
داده‌شده  در کلاس اول یا کلاس دوم باشند را می‌دهد. احتمال این که در کلاس اول باشد، برابر است با
$1/(1+\exp(-r(\x,\uu))$.
فرض می‌کنیم که 
$h_i$
ها و 
$\psi_i$
ها ویژگی
\lr{Universal Approximation Capacity}
را داشته باشند. در ادامه به بیان دقیق قضیه‌ی جدایی‌پذیری  می‌پردازیم:

\begin{theorem}
	فرض کنید:
\begin{enumerate}
	\item 
فرض کنید داده از مدل غیرخطی
$\x = \f(\s)$
تولید می‌شود و رابطه‌ی 
\eqref{genprior}
برقرار است.
\item
به ازای هر 
$\uu$
توابع 
$q_i$
به عنوان تابعی از 
$s_i$
به اندازه‌ی کافی هموار هستند.
\item
برای هر بردار 
$\y \in \R^n$،
تعداد 
$2n+1$
مقدار برای 
$\uu$
وجود دارد که با 
$\uu_i$
 نمایش داده‌شده اند. این بردار‌ها به گونه‌ای هستند که
 $2n$
 بردار
 \begin{multline}
 \left(\w(\y,\uu_1)-\w(\y,\uu_0)),(\w(\y,\uu_2)-\w(\y,\uu_0)),\right. \\ \left. ...,(\w(\y,\uu_{\ncond})-\w(\y,\uu_0)\right)
 \end{multline}
که عضو
 $\R^{2n}$
 هستند و در آن‌ها 
 \begin{multline} \label{wdef}
 \w(\y,\uu)=\left(\frac{\partial q_1(y_1,\uu)}{ \partial y_1},\ldots,\frac{\partial q_n(y_n,\uu)}{ \partial y_n},  \right. \\ \frac{\partial^2 q_1(y_1,\uu)}{ \partial y_1^2},   \left. \ldots,\frac{\partial^2 q_n(y_n,\uu)}{ \partial y_n^2}\right)
 \end{multline}
 است، مستقل خطی باشند.
\item 
یک رگرسون لاجیستیک
با  
\lr{Universal Approximation Capacity}
را آموزش می‌دهیم تا بین 
$\xaug$
و
$\xaug^*$
تمایز قائل شود.
\item
فرض می‌کنیم توابع 
$
r(\x,\uu)= \sum_{i=1}^n \psi_i(h_i(\x),\uu)
$
هموار و وارون‌پذیر بوده و وارون آن‌ها نیز هموار باشد.
\end{enumerate}
در این صورت، در حدی که بیشمار داده در اختیار داریم، 
 $h_i(\x)$
برابر منابع مستقل با حداکثر اعمال یک تابع  وارون‌پذیر هستند.
\label{newthm}
\end{theorem}
\begin{proof}
در این‌جا، از آوردن اثبات این قضیه چشم‌پوشی کرده‌ایم. اثبات این قضیه در بخش اول فصل
\lr{Supplementary Material}
مقاله‌ی 
\cite{newHyva}
آمده است.
\end{proof}

\begin{definition}
متغیر تصادفی 
$s_i$
به شرط 
$\uu$
دارای توزیع نمایی شرطی از مرتبه‌ی 
$k$
است، اگر 
\lr{pdf}
آن به ازای هر 
$\uu$
 قابل نوشتن به فرم زیر باشد:
\begin{equation} \label{expdef}
p(s_i|\uu)=\frac{Q_i(s_i)}{Z_i(\uu)}\exp[\sum_{j=1}^k \suffstat_{ij}(s_i)\lambda_{ij}(\uu)]
\end{equation}
در این معادله،
$Q_i$،
$\lambda_{ij}$،
$\tilde{q}_{ij}$
و 
$Z_i$
توابعی اسکالر هستند.
\end{definition}
\begin{example}
در هر زمان داده‌شده، مقدار یک فرآیند تصادفی گوسی با 
$\uu$
ای برابر گذشته‌ی این فرآیند در این زمان، یک متغیر تصادفی نمایی شرطی با پارامتر 
$k = 1$
است.
\end{example}

\begin{theorem}
اگر منابع
$s_i$
به شرط 
$\uu$
 دارای توزیع شرطی نمایی با مرتبه‌ی 
$k$
 باشند. در صورتی که  $k=1$ باشد، شرط سوم قضیه‌ی 
\eqref{newthm}
نمی‌تواند برقرار باشد.
\end{theorem}
\begin{proof}
 اثبات این قضیه در بخش دوم فصل
\lr{Supplementary Material}
مقاله‌ی 
\cite{newHyva}
آمده است.
\end{proof}
\begin{note}
قضیه‌ی فوق، مشابه جدایی‌ناپذیر بودن متغیر‌های تصادفی گوسی در حالت خطی است.

\section{استفاده از مشتق منابع}
فرض کنید منابع 
$\s$
به صورت غیر خطی با یکدیگر مخلوط‌ شده باشند:
\begin{equation}
\x(t) = \f(\s(t))
\end{equation}
در این حالت، با وجود غیرخطی بودن تابع 
$\f$،
مشتق منابع در هر لحظه‌ به صورت خطی با یکدیگر مخلوط شده‌اند و مشتق سیگنال‌های 
$\x$
را تولید می‌کنند:
\begin{equation}
\frac{dx_i}{dt} = \sum_{j = 1}^{n} \frac{\partial f_i}{\partial s_j} \frac{ds_j}{dt}
\implies
\nabla\x = \mathbf{J}_{\f; t}(\s)\nabla\s
\end{equation}
که در آن، ماتریس 
$ \mathbf{J}_{\f; t}$
برابر ژاکوبین تبدیل 
$\f$
در نقطه‌ی 
$\s(t)$
است. 
\begin{figure}[h!]
	\centering
	\includegraphics[scale=0.4]{ehsan.png}
	\caption{بهره‌بردن از مشتق برای جداسازی مخلوط‌های غیرخطی}
	\label{ehsan}
\end{figure}

این ایده برای اولین بار در مقاله‌ی 
\cite{ehsandoust}
مطرح شده است. ممکن است به نظر برسد که بر مبنای این ایده‌، و با توجه به قضیه‌ی جدایی‌پذیری مخلوط‌های خطی، به سادگی می‌توان قضیه‌ای برای جدایی‌پذیری مخلوط‌های غیرخطی مطرح کرد، اما این‌گونه نیست زیرا در این‌جا، ماتریس
$ \mathbf{J}_{\f; t}$
تابعی از زمان است و ثابت نیست. با این وجود، می‌توان بر مبنای این ایده و با فرض اینکه تغییرات ماتریس ژاکوبین در زمان بسیار کند است، الگوریتمی وفقی برای جداسازی کور مخلوط غیرخطی ارائه داد که تغییرات کند ماتریس ژاکوبین را دنبال کند. این روند در 
\cite{ehsandoust}
پی گرفته شده و الگوریتمی برای جداسازی در حالت تغییر کند ماتریس ژاکوبین ارائه گشته است.  


\end{note}